\documentclass[ ../main.tex]{subfiles}
\providecommand{\mainx}{..}
\begin{document}
\appendix
\appendixpage
\addappheadtotoc
\section{Probability mass of random bit length}
\begin{algorithm}
    \caption{Bit length sampler of the cryptographic perfect hash function}
    \label{alg:sampler}
    \SetKwProg{func}{function}{}{}
    \KwIn
    {
        The number $m$ is the cardinality of the perfect hash set and
        the number $r$ is the load factor.
    }
    \KwOut
    {
        The \emph{minimum bit length} $n$ conditioned on a random set of cardinality $m$ and a \gls{gls-loadfactor} $r$.
    }
    \func{$\sampler(m, r)$}
    {
        \For{$i \gets 1$ \KwTo $m$}
        {
            $x \gets \text{randomly draw a bit string from $\BitSet^{*}$ without replacement}$\;
            $\Set{S} \gets \Set{S} \cup \{x\}$\;
        }    
        $(b_n,N) \gets \makehash\!\left(\Set{S}, r\right)$\;
        \tcp{The integer $N$ must also be coded, which we assume has a constant bit length.}
        \Return $n + \mathcal{O}(1)$\;
    }
\end{algorithm}
The bit length of the perfect hash function found by \cref{alg:ph} has an uncertain value with respect to (random) sets and therefore we may model it as a random variable as given by the following definition.
\begin{definition}
\label{def:N_rv}
The perfect hash function generated by \cref{alg:ph} has a random bit length given by
\begin{equation}
    \RV{N} = \sampler(m,r)
\end{equation}
where $m$ is the cardinality of the random set and $r$ is the load factor.
\end{definition}
Note that if the distribution of sets is not uniformly distributed as given by 
\sampler, then better algorithms than \cref{alg:ph} are possible, i.e., the 
lower-bounds given by \cref{??} are with respect to ``uniformly distributed'' 
random sets and smaller lower-bounds are possible if the random sets are 
non-uniformly distributed.

The probability mass of the random bit length $\RV{N}$ is given by the following theorem.
\begin{theorem}
\label{thm:N_pmf}
The random bit length $\RV{N}$ has a probability mass function given by
\begin{equation}
\label{eq:N_pmf}
    \pmf_{\RV{N}}(n \Given m, r) = q^{2^n-1}\left(1 - q^{2^n}\right)\,.
\end{equation}
where $m$ is the cardinality of the random set, $r$ is the load factor, and $q = 1 - p(m,r)$, 
\end{theorem}
\begin{proof}
Each iteration of the loop in \cref{alg:ph} has a collision test which is Bernoulli distributed with a probability of success $p(m,r)$, where success denotes no collision occurred. We are interested in the random length $\RV{N}$ of the bit string when this outcome occurs.

For the random variable $\RV{N}$ to realize a value $n$, every bit string smaller than length $n$ must fail and a bit string of length $n$ must succeed. There are $2^n-1$ bit strings smaller than length $n$ and each one fails with probability $q$, and so by the product rule the probability that they all fail is given by
\begin{equation}
\label{eq:N_pmf_1}
    q^{2^n - 1}\,.
\end{equation}

Given that every bit string smaller than length $n$ fails, what is the probability that every bit string of length $n$ fails? There are $2^n$ bit strings of length $n$, each of which fails with probability $q$ as before and thus by the product rule the probability that they all fail is $q^{2^n}$, whose complement, the probability that not all bit strings of length $n$ fail, is given by
\begin{equation}
\label{eq:N_pmf_2}
    1 - q^{2^n}\,.
\end{equation}
By the product rule, the probability that every bit string smaller than length $n$ fails and a bit string of length $n$ succeeds is given by the product of (\ref{eq:N_pmf_1}) and (\ref{eq:N_pmf_2}),
\begin{equation}
\label{eq:prob_is_mass}
    q^{2^n-1}\left(1 - q^{2^n}\right)\,.
\end{equation}
For \cref{eq:prob_is_mass} to be a probability mass function, two conditions must be met. First, its range must be a subset of $[0,1]$. Second, the summation over its domain must be $1$.

The first case is trivially shown by the observation that $q$ is a positive number between $0$ and $1$ and therefore any non-negative power of $q$ is positive number between $0$ and $1$.

The second case is shown by calculating the infinite series
\begin{align}
    S &= \sum_{n=0}^{\infty} q^{2^n-1}\left(1-q^{2^n}\right)\\
      &= \sum_{n=0}^{\infty} q^{2^n-1} - q^{2^{n+1}-1}\,.
\end{align}
Explicitly evaluating this series for the first $4$ terms reveals a telescoping sum given by
\begin{equation}
    S = (1 - q) + (q - q^3) + (q^3 - q^7) + (q^7 - q^{15}) + \cdots\,,
\end{equation}
where everything cancels except $1$.
\end{proof}

In \cref{fig:pmf_fN}, we plot the probability mass function of $\RV{N}$ conditioned on $m=25$ and several different load factors. We see that the probability mass is peaked and unimodal and tends to nearly zero everywhere except over a concentrated interval around its expected value.
\begin{figure}
\centering
\caption{Probability mass function of\\random bit length $\RV{N}$}
\label{fig:pmf_fN}
\includegraphics{plotfn}
\end{figure}

The expected size was already computed, but the probability mass function contains everything there is to know about the distribution of $\RV{N}$, not just the \emph{expected} value. However, for illustration, we show how the expected value may be computed.
\begin{theorem}
\label{thm:open_expect}
The expected bit length of $\RV{N}$ conditioned on random sets of cardinality $m$ and a load factor $r$ is given by
\begin{equation}
\label{eq:alt_expectN}
    \Expect{\RV{N}} = \sum_{j=1}^{\infty} q^{2^j-1}\,,
\end{equation}
where $q = 1 - p(m,r)$.
\end{theorem}
\begin{proof}
The expectation of $\RV{N}$ is given by
\begin{align}
    &\sum_{j=0}^{\infty} j q^{2^j-1}\left(1 - q^{2^j}\right)\\
    &\qquad=\sum_{j=0}^{\infty} j \left(q^{2^j-1} - q^{2^{j+1}-1}\right)\,.
\end{align}
Explicitly evaluating this series for the first $4$ terms reveals a converging sum given by
\begin{align}
    &0 + (q - q^3) + 2(q^3 - q^7) + 3(q^7 - q^{15}) + \cdots\\
    &\qquad= q + q^3 + q^7 + q^{15} + \cdots\,.
\end{align}
\end{proof}
When we numerically evaluate \cref{eq:alt_expectN} for various $m$ and $r$, the results are in agreement with \cref{eq:code_size}.
\end{document}